\documentclass[a4paper,11pt]{article}

\usepackage[T1]{fontenc}
\usepackage{geometry}
\usepackage{hyperref}
\usepackage{parskip}
\usepackage{cfr-lm}
\usepackage{fontawesome5}

\geometry{left=2.5cm, right=2.5cm, top=2.5cm, bottom=2.5cm}

\hypersetup{
    colorlinks=true,
    urlcolor=blue!70!black
}

\pagestyle{empty}

\newcommand{\headerfontiii}{\fontfamily{ppl}\selectfont}

\begin{document}
\headerfontiii

\begin{flushleft}
\textbf{Loghman Samani} \\
Ludwigsburg, Baden-W\"urttemberg, Germany \\
\href{mailto:samaniloqman91@gmail.com}{samaniloqman91@gmail.com} \\
\href{https://loqmansamani.github.io}{loqmansamani.github.io}
\end{flushleft}

\vspace{4mm}

\begin{flushleft}
February 25, 2026
\end{flushleft}

\vspace{4mm}

\begin{flushleft}
Helmholtz Munich \\
Institute of Computational Biology \\
Ingolst\"adter Landstra\ss e 1 \\
85764 Neuherberg, Germany
\end{flushleft}

\vspace{4mm}

\textbf{Application for PhD Candidate (f/m/x) -- AI/ML Drug Discovery for Brain Diseases (102846)}

\vspace{4mm}

Dear Hiring Committee,

I would like to apply for the PhD position in AI/ML Drug Discovery for Brain Diseases at the Institute of Computational Biology, Helmholtz Munich.

I am finishing my Master's in Technical Biology at the University of Stuttgart, where my thesis focused on designing gene regulatory networks through a combination of evolutionary algorithms and gradient-based optimization over PDE-based reaction-diffusion models. Before that, I studied Cell and Molecular Biology at the University of Kurdistan. So I have been moving between biology and computation for a while now, and this position feels like a natural fit for where I want to go next.

Most of my recent work has gone into building open-source tools for deep learning and computational biology. TorchDiff, my PyTorch library for diffusion models, has been downloaded around 5,000 times and supports DDPM, DDIM, SDE-based, and latent diffusion frameworks. I also built CppNet, a deep learning library written entirely in C++, and BioStoch, a library for stochastic and deterministic simulation of biological systems. Alongside these, I have implemented various architectures from scratch, including CNNs, RNNs, Transformers, and a GPT-2 model. I work primarily in Python and PyTorch, and I am comfortable with C++ and numerical methods.

What draws me to this position is the chance to apply graph neural networks, deep learning, and optimal transport methods to real patient data from neurological diseases. Working with spatial transcriptomics and single-cell data to understand disease mechanisms and identify drug targets is a genuinely exciting problem. I have spent enough time in both molecular biology and machine learning to be able to engage with the biological questions and the computational challenges of a project like this, and I would be glad to contribute to the collaborative effort with the Fugger lab.

I would be happy to discuss my background in more detail. Thank you for considering my application.

\vspace{6mm}

Best regards, \\
Loghman Samani

\end{document}
